\subsection{李代数中的括号}

对于 \(\alpha \in \mathbf{N}^n\)，令 \(e_{\alpha}\) 为原点处的点分布 \(\frac{1}{\alpha!} \frac{\partial^{\alpha}}{\partial x^{\alpha}}\)。特别地，\(e_k = e_{\varepsilon_k} = \frac{\partial}{\partial x_k}\)。这些 \(e_{\alpha}\) 构成向量空间 \(\mathrm{U}(G)\) 的一组基。若 \(f\) 是 G 中 0 的一个开邻域上的解析函数，且 \(f(x) = \sum_{\alpha} \lambda_{\alpha} x^{\alpha}\) 是它在原点附近的幂级数展开，则

特别地，

\[
\langle e _ {\alpha}, f \rangle = \lambda_ {\alpha}.
\]

\[
\langle e _ {\alpha}, x ^ {\gamma} \rangle = \delta_ {\alpha \gamma}.
\]

因此

\[
\begin{array}{r l} \langle e _ {\alpha} * e _ {\beta}, x ^ {\gamma} \rangle & = \langle e _ {\alpha} \otimes e _ {\beta}, (x. y) ^ {\gamma} \rangle \\ & = \langle e _ {\alpha} \otimes e _ {\beta}, \sum_ {\alpha^ {\prime}, \beta^ {\prime}} c _ {\alpha^ {\prime} \beta^ {\prime} \gamma} x ^ {\alpha^ {\prime}} y ^ {\beta^ {\prime}} \rangle \\ & = \sum_ {\alpha^ {\prime}, \beta^ {\prime}} c _ {\alpha^ {\prime} \beta^ {\prime} \gamma} \langle e _ {\alpha}, x ^ {\alpha^ {\prime}} \rangle \langle e _ {\beta}, y ^ {\beta^ {\prime}} \rangle = c _ {\alpha \beta \gamma} \end{array}
\]

从而

\[
e _ {\alpha} * e _ {\beta} = \sum_ {\gamma} c _ {\alpha \beta \gamma} e _ {\gamma}.\tag{11}
\]

（公式 (10) 表示 U(G) 上的结合律。）

特别地，由于 L(G) 在括号运算下稳定，

\[
[ e _ {i}, e _ {j} ] = \sum_ {k} (c _ {i j k} - c _ {j i k}) e _ {k}.\tag{12}
\]

因此，\(\mathbf{L}(\mathbf{G})\) 相对于基 \((e_{1},\ldots,e_{n})\) 的结构常数为 \(c_{ijk}-c_{jik}\)。换言之，将 \(\mathbf{L}(\mathbf{G})\) 规范地识别为 \(K^{n}\)，得到：

\[
[ x, y ] = \mathrm{B} (x, y) - \mathrm{B} (y, x).\tag{13}
\]

命题 1.

\begin{enumerate}
\def\labelenumi{(\roman{enumi})}
\tightlist
\item
\end{enumerate}

\[
x ^ {[ - 1 ]} \equiv - x + \mathrm{B} (x, x) \mod \deg 3\tag{i}
\]

\[
x. y. x ^ {[ - 1 ]} \equiv y + [ x, y ] \quad \text {   mod   } \deg 3\tag{ii}
\]

\[
y ^ {[ - 1 ]}. x. y \equiv x + [ x, y ] \quad \text {   mod   } \deg 3\tag{iii}
\]

\[
x ^ {[ - 1 ]}. y ^ {[ - 1 ]}. x. y \equiv [ x, y ] \quad \text {   mod   } \deg 3\tag{iv}
\]

\[
x. y. x ^ {[ - 1 ]}. y ^ {[ - 1 ]} \equiv [ x, y ] \quad \text {   mod   } \deg 3\tag{v}.
\]

（在 (i) 中，\(x^{[-1]}\) 表示函数 \(x \mapsto x^{[-1]}\) 在原点附近的幂级数展开；其他公式可以类似解释。）

令 \(g_{1}\) 和 \(g_{2}\) 为 \(x^{[-1]}\) 的 1 次和 2 次齐次分量。于是

\[
\begin{array}{r l} 0 = x. x ^ {[ - 1 ]} \\ \equiv x + g _ {1} (x) + \mathrm{B} (x, g _ {1} (x)) & \text {   mod   } \deg 2 \quad \text {(by (6))} \\ \equiv x + g _ {1} (x) & \text {   mod   } \deg 2 \end{array}
\]

从而 \(g_{1}(x) = -x\)。于是

\[
\begin{array}{r l} 0 & = x. x ^ {[ - 1 ]} \\ & \equiv x + (- x + g _ {2} (x)) + B (x, - x + g _ {2} (x)) \mod \deg 3 \\ & \equiv g _ {2} (x) - B (x, x) \mod \deg 3 \end{array}
\]

从而 \(g_{2}(x) = \mathrm{B}(x, x)\)。这证明了 (i)。然后利用 (i)，

\[
\begin{array}{l l} x. y. x ^ {[ - 1 ]} \equiv (x + y + \mathrm{B} (x, y)) \cdot (- x + \mathrm{B} (x, x)) & \text {mod} \deg 3 \\ \equiv x + y + \mathrm{B} (x, y) + (- x + \mathrm{B} (x, x)) + \mathrm{B} (x + y, - x) & \text {mod} \deg 3 \\ \equiv y + \mathrm{B} (x, y) - \mathrm{B} (y, x) & \text {mod} \deg 3 \\ \equiv y + [ x, y ] \qquad \qquad \qquad \qquad \qquad \qquad \qquad \qquad \qquad \qquad \qquad \qquad \qquad \qquad \qquad \qquad \qquad \qquad \qquad \qquad \qquad \qquad \qquad \qquad \qquad \qquad \qquad \qquad \qquad \qquad \qquad \qquad \qquad \qquad \qquad \qquad \qquad \\ & (\text {by (13)}) \end{array}
\]

从而得到 (ii)。(iii) 的证明类似。结合 (i) 和 (iii)，得到

\[
\begin{array}{r l r} x ^ {[ - 1 ]}. y ^ {[ - 1 ]}. x. y & \equiv (- x + \mathrm{B} (x, x)). (x + [ x, y ]) & \bmod \deg 3 \\ & \equiv - x + \mathrm{B} (x, x) + x + [ x, y ] + \mathrm{B} (- x, x) & \bmod \deg 3 \\ & \equiv [ x, y ] & \bmod \deg 3 \end{array}
\]

从而得到 (iv)。(v) 的证明类似。

